\documentclass[11pt]{article}
\pagestyle{plain}

% \pgfplotsset{compat=1.18}
\usepackage[margin=1in]{geometry}
\usepackage{amsmath,amsfonts,amssymb,amsthm}
\usepackage[ruled,vlined,linesnumbered,commentsnumbered]{algorithm2e}
\SetKwInput{KwIn}{Input}
\SetKwInput{KwOut}{Output}
\SetKwInput{KwGlobal}{Global variable}
\SetKwInput{KwOracle}{Oracle}
\SetKwComment{Comment}{// }{}
\usepackage{graphicx}
\usepackage{enumerate}
\usepackage{bbm}
\usepackage{verbatim}
\usepackage{hyperref,color}
\usepackage[capitalize,nameinlink]{cleveref}
\usepackage[dvipsnames]{xcolor}
\hypersetup{
	colorlinks=true,
	pdfpagemode=UseNone,
	citecolor=OliveGreen,
	linkcolor=NavyBlue,
	urlcolor=Magenta,
	pdfstartview=FitW
}
\usepackage{appendix}
\usepackage{makecell}
\usepackage{tablefootnote}
\crefname{appsec}{Appendix}{Appendices}
\usepackage{tikz}
\usepackage{pgfplots}
\usepackage{xifthen}
\usepackage{subcaption}
\usepackage{hhline}

\theoremstyle{plain}
\newtheorem{theorem}{Theorem}[section]
\newtheorem{proposition}[theorem]{Proposition}
\newtheorem{lemma}[theorem]{Lemma}
\newtheorem{corollary}[theorem]{Corollary}
\newtheorem{conjecture}[theorem]{Conjecture}
\newtheorem{problem}[theorem]{Problem}
\newtheorem{claim}[theorem]{Claim}
\newtheorem{fact}[theorem]{Fact}

\theoremstyle{definition}
\newtheorem{definition}[theorem]{Definition}
\newtheorem{example}[theorem]{Example}
\newtheorem{observation}[theorem]{Observation}
\newtheorem{assumption}[theorem]{Assumption}
\newtheorem*{assumption*}{Assumption}
\newtheorem{condition}[theorem]{Condition}

\theoremstyle{remark}
\newtheorem{remark}[theorem]{Remark}

\crefname{lemma}{Lemma}{Lemmas}
\crefname{theorem}{Theorem}{Theorems}
\crefname{definition}{Definition}{Definitions}
\crefname{fact}{Fact}{Facts}
\crefname{claim}{Claim}{Claims}
\crefname{proposition}{Proposition}{Propositions}


\newcommand{\dif}{\,\mathrm{d}}
%\newcommand{\E}{\mathbb{E}}
%\newcommand{\Var}{\mathrm{Var}}
\newcommand{\Cov}{\mathrm{Cov}}
\newcommand{\MI}{\mathrm{I}}
%\newcommand{\Ent}{\mathrm{Ent}}
\newcommand{\ind}{\mathbbm{1}}
\newcommand{\allone}{\mathbf{1}}
\newcommand{\tr}{\mathrm{tr}}
\newcommand{\wt}{\mathrm{weight}}
\DeclareMathOperator*{\argmax}{arg\,max}
\newcommand{\norm}[1]{\left\lVert #1 \right\rVert}
\newcommand{\TV}[2]{\left\|#1 - #2\right\|_{\mathrm{TV}}}
\newcommand{\ceil}[1]{\left\lceil #1 \right\rceil}
\newcommand{\floor}[1]{\left\lfloor #1 \right\rfloor}
\newcommand{\ip}[2]{\left\langle #1 , #2 \right\rangle}
\newcommand{\grad}{\nabla}
\newcommand{\hessian}{\nabla^2}
\newcommand{\trans}{\intercal}
\newcommand{\diag}{\mathrm{diag}}
\newcommand{\poly}{\mathrm{poly}}
\newcommand{\dist}{\mathrm{dist}}
\newcommand{\sgn}{\mathrm{sgn}}
\newcommand{\fpras}{\mathsf{FPRAS}}
\newcommand{\fpaus}{\mathsf{FPAUS}}
\newcommand{\fptas}{\mathsf{FPTAS}}

\newcommand{\eps}{\varepsilon}

\newcommand{\N}{\mathbb{N}}
\newcommand{\Z}{\mathbb{Z}}
\newcommand{\Q}{\mathbb{Q}}
\newcommand{\R}{\mathbb{R}}
\newcommand{\C}{\mathbb{C}}

\renewcommand{\AA}{\mathcal{A}}
\newcommand{\BB}{\mathcal{B}}
\newcommand{\CC}{\mathcal{C}}
\newcommand{\DD}{\mathcal{D}}
\newcommand{\EE}{\mathcal{E}}
\newcommand{\FF}{\mathcal{F}}
\newcommand{\GG}{\mathcal{G}}
\newcommand{\HH}{\mathcal{H}}
\newcommand{\II}{\mathcal{I}}
\newcommand{\JJ}{\mathcal{J}}
\newcommand{\KK}{\mathcal{K}}
\newcommand{\LL}{\mathcal{L}}
\newcommand{\MM}{\mathcal{M}}
\newcommand{\NN}{\mathcal{N}}
\newcommand{\OO}{\mathcal{O}}
\newcommand{\PP}{\mathcal{P}}
\newcommand{\QQ}{\mathcal{Q}}
\newcommand{\RR}{\mathcal{R}}
\renewcommand{\SS}{\mathcal{S}}
\newcommand{\TT}{\mathcal{T}}
\newcommand{\UU}{\mathcal{U}}
\newcommand{\VV}{\mathcal{V}}
\newcommand{\WW}{\mathcal{W}}
\newcommand{\XX}{\mathcal{X}}
\newcommand{\YY}{\mathcal{Y}}
\newcommand{\ZZ}{\mathcal{Z}}

\newcommand{\KL}[2]{D_{\mathrm{KL}} \left( #1 \,\Vert\, #2 \right)}
\newcommand{\KLbig}[2]{D_{\mathrm{KL}} \left( #1 \,\big\Vert\, #2 \right)}
\newcommand{\KLBig}[2]{D_{\mathrm{KL}} \left( #1 \,\Big\Vert\, #2 \right)}

\newcommand{\SKL}[2]{D_{\mathrm{SKL}} \left( #1 \,\Vert\, #2 \right)}
\newcommand{\SKLbig}[2]{D_{\mathrm{SKL}} \left( #1 \,\big\Vert\, #2 \right)}
\newcommand{\SKLBig}[2]{D_{\mathrm{SKL}} \left( #1 \,\Big\Vert\, #2 \right)}

\newcommand{\qandq}{\quad\text{and}\quad}
\newcommand{\supp}{\mathrm{supp}}

\newcommand{\DKL}[2]{\-D_{\-{KL}}\left(#1 \parallel #2\right)}
\newcommand{\DTV}[2]{\-D_{\mathrm{TV}}\left({#1},{#2}\right)}
% \newcommand{\dist}{\mathrm{dist}}
\newcommand{\e}{\mathrm{e}}
\renewcommand{\epsilon}{\varepsilon}
\renewcommand{\emptyset}{\varnothing}
% \newcommand{\norm}[1]{\left\Vert#1\right\Vert}
\newcommand{\set}[1]{\left\{#1\right\}}
\newcommand{\tuple}[1]{\left(#1\right)} 
% \newcommand{\eps}{\varepsilon}
\newcommand{\inner}[2]{\left\langle #1,#2\right\rangle}
\newcommand{\tp}{\tuple}
\newcommand{\ol}{\overline}
\newcommand{\abs}[1]{\left\vert#1\right\vert}
\newcommand{\ctp}[1]{\left\lceil#1\right\rceil}
\newcommand{\ftp}[1]{\left\lfloor#1\right\rfloor}
\newcommand{\Ind}[1]{\-{Ind}_{#1}}
\newcommand{\todo}[1]{{\color{red} TODO: #1}}

\newcommand{\Cor}{\Psi^{\mathrm{sym}}}
\newcommand{\cor}{\Psi^{\-{cor}}}
\newcommand{\Inf}{\Psi}
% \newcommand{\diag}{\-{diag}}

\def\*#1{\boldsymbol{#1}} % Use \*A for \mathbf{A}
\def\+#1{\mathcal{#1}} % Use \+A for \mathcal{A}
\def\-#1{\mathrm{#1}} % Use \-A for \mathrm{A}
\def\=#1{\mathbb{#1}} % Use \^A for \mathbb{A}
\def\!#1{\mathfrak{#1}} % Use \!A for \mathfrak{A}

\newcommand*\diff{\mathop{}\!\mathrm{d}}

% probability
% \DeclareMathOperator*{\oPr}{\mathbf{Pr}}
\def\oPr{\mathop{\mathrm{Pr}}}
\renewcommand{\Pr}[2][]{ \ifthenelse{\isempty{#1}}
  {\oPr\left[#2\right]}
  {\oPr_{#1}\left[#2\right]} } % Use \Pr[a]{b} for \mathbf{Pr}_a[b], \Pr{b} for  \mathbf{Pr}[b]

% expectation: relies on the package "dsfont"
% \DeclareMathOperator*{\oE}{\mathbb{E}}
\def\oE{\mathop{\mathbb{E}}}
\newcommand{\E}[2][]{ \ifthenelse{\isempty{#1}}
  {\oE\left[#2\right]}
  {\oE_{#1}\left[#2\right]} }

% variance
%\DeclareMathOperator*{\oVar}{\mathbf{Var}}
\def\oVar{\mathrm{Var}}
\newcommand{\Var}[2][]{ \ifthenelse{\isempty{#1}}
  {\oVar\left[#2\right]}
  {\oVar_{#1}\left[#2\right]} }

% entropy
% \DeclareMathOperator*{\oEnt}{\mathbf{Ent}}
\def\oEnt{\mathrm{Ent}}
\newcommand{\Ent}[2][]{ \ifthenelse{\isempty{#1}}
  {\oEnt\left[#2\right]}
  {\oEnt_{#1}\left[#2\right]} }

\newcommand{\PhiEnt}[2][]{ \ifthenelse{\isempty{#1}}
  {\oEnt^\phi\left[#2\right]}
  {\oEnt^\phi_{#1}\left[#2\right]} }


\newenvironment{Exampleparagraph}[1][]
{
    \par 
    \noindent 
    \textbf{\color{blue} #1}
    \par 
    \vspace{0.5em} 
    \itshape 
}
{
    \par 
    \vspace{1em} 
}

%\input{local-macro}

\newcommand{\RED}[1]{\textcolor{red}{#1}}
\newcommand{\zc}[1]{\textcolor{red}{ZC: #1}}

\newcommand{\bern}[1]{\mathrm{Bern}\left(#1\right)}



\title{Local Bobkov inequality and Eldan-Gross inequality}
\date{\today}

\pgfplotsset{compat=1.18} 
\begin{document}

\maketitle

\section*{Disclaimer}
Calculations in this note are done by GPT and codex.

\section{Local Bobkov inequality for general measure}

Local Bobkov inequality:
\begin{align*}
    I(P_t f) \le P_t \tp{\sqrt{I(f)^2 + M (1-\e^{-\rho t}) \abs{\grad f}^2}}.
\end{align*}

Let $\mu$ be a measure on $\{-1,+1\}^n$ and $Lf (x) = \sum_{i \in [n]} p_i(x) \tp{f(x^{\oplus i}) - f(x)}$ be the generator, where $p_i(x) = p(x,x^{\oplus i})$ denotes the transition rate from $x$ to $x^{\oplus i}$.
Define the gradient $\nabla f(x)$ as follows:
\begin{align*}
    \nabla f(x) = \tp{\partial_i f(x)}_{i \in [n]} = \tp{p_i(x) (f(x^{\oplus i}) - f(x))}_{i \in [n]}.
\end{align*}
Let $u_s=f_{t-s}=P_{t-s}f$ and $a_s=\sqrt{M(1-\e^{-\rho s})}$. Define
\begin{align*}
    F_s(x) = \sqrt{I(u_s(x))^2+a_s^2\abs{\nabla u_s(x)}^2},
    \qquad
    V_s(x) = \tp{I(u_s(x)), a_s\partial_1u_s(x), \ldots, a_s\partial_nu_s(x)}.
\end{align*}
By a similar calculation, we have
\begin{align*}
    L F_s(x) = \frac{\inner{V_s}{L V_s}}{\abs{V_s}}(x) + \sum_{i \in [n]} p_i(x) \underbrace{\tp{\abs{V_s}(x^{\oplus i}) - \frac{\inner{V_s(x)}{V_s(x^{\oplus i})}}{\abs{V_s}(x)}}}_{:=D_i(x)}.
\end{align*}
Therefore, we have
\begin{align*}
    (\partial_s + L) F_s(x) = \frac{\inner{V_s}{(\partial_s + L) V_s}}{\abs{V_s}(x)} + \sum_{i \in [n]} p_i(x) D_i(x).
\end{align*}
The second term can be bounded as follows. Since
\begin{align*}
    (V_s(x^{\oplus i}))_i
    = -\frac{p_i(x^{\oplus i})}{p_i(x)}(V_s(x))_i,
\end{align*}
we have
\begin{align*}
    D_i(x)
    &\ge \frac{2 p_i(x^{\oplus i})}{p_i(x)} \frac{(V_{s}^2(x))_i}{\abs{V_s}(x)},\\
    p_i(x)D_i(x)
    &\ge 2 p_i(x^{\oplus i}) \frac{(V_{s}^2(x))_i}{\abs{V_s}(x)}.
\end{align*}
By a straightforward calculation, the first term satisfies
\begin{align*}
    ((\partial_s + L) V_s)_0(x)
    &= L I(u_s)(x)-I'(u_s(x))Lu_s(x),\\
    ((\partial_s + L) V_s)_i(x)
    &= \frac{M\rho\e^{-\rho s}}{2a_s}\partial_i u_s(x)
    +a_s\tp{L\partial_i-\partial_i L}u_s(x)\\
    &= \frac{M\rho\e^{-\rho s}}{2a_s}\partial_i u_s(x)
    -a_s[\partial_i,L]u_s(x)\\
    &= \frac{M\rho\e^{-\rho s}}{2a_s}\partial_i u_s(x)
    -a_s\sum_{j \in [n]}[\partial_i,\partial_j]u_s(x),
    \qquad i \in [n],
\end{align*}
where $[A,B]=AB-BA$. Writing $\delta_i u_s(x)=u_s(x^{\oplus i})-u_s(x)$, the zeroth component is bounded by
\begin{align*}
    I(u_s(x))((\partial_s+L)V_s)_0(x)
    &= I(u_s(x))\sum_{i \in [n]}p_i(x)\tp{I(u_s(x^{\oplus i}))-I(u_s(x))-I'(u_s(x))\delta_i u_s(x)}\\
    &\ge -\sum_{i \in [n]}p_i(x)\tp{\delta_i u_s(x)}^2,
\end{align*}
where we use
\begin{align*}
    I(a)\tp{I(b)-I(a)-I'(a)(b-a)}\ge -(b-a)^2.
\end{align*}
Combining this with the bounds on $D_i(x)$, it holds
\begin{align*}
    (\partial_s+L)F_s(x)
    &\ge \frac{1}{\abs{V_s(x)}}\Bigg[
    -\sum_{i \in [n]}p_i(x)\tp{\delta_i u_s(x)}^2
    +\frac{M\rho\e^{-\rho s}}{2}\sum_{i \in [n]}\tp{\partial_i u_s(x)}^2\\
    &\qquad
    +M(1-\e^{-\rho s})\sum_{i \in [n]}\partial_i u_s(x)
    \tp{L\partial_i-\partial_i L}u_s(x)\\
    &\qquad
    +2M(1-\e^{-\rho s})\sum_{i \in [n]}p_i(x^{\oplus i})
    \tp{\partial_i u_s(x)}^2
    \Bigg].
\end{align*}
Equivalently, on edges with $p_i(x)>0$,
\begin{align*}
    (\partial_s+L)F_s(x)
    &\ge \frac{1}{\abs{V_s(x)}}\Bigg[
    \sum_{i \in [n]}\tp{\frac{M\rho\e^{-\rho s}}{2}
    +2M(1-\e^{-\rho s})p_i(x^{\oplus i})
    -\frac{1}{p_i(x)}}\tp{\partial_i u_s(x)}^2\\
    &\qquad
    +M(1-\e^{-\rho s})\sum_{i \in [n]}\partial_i u_s(x)
    \tp{L\partial_i-\partial_i L}u_s(x)
    \Bigg].
\end{align*}



\paragraph{Special choice of rates.}
Suppose now that
\begin{align*}
    p_i(x)=\sqrt{\frac{\mu(x^{\oplus i})}{\mu(x)}}.
\end{align*}
Then the products along the two paths of any square agree:
\begin{align*}
    p_j(x)p_i(x^{\oplus j})
    =
    \sqrt{\frac{\mu(x^{\oplus j})}{\mu(x)}}
    \sqrt{\frac{\mu(x^{\oplus i\oplus j})}{\mu(x^{\oplus j})}}
    =
    p_i(x)p_j(x^{\oplus i}).
\end{align*}
Thus the second-order term in the commutator cancels. Indeed, for $i \neq j$,
\begin{align*}
    [\partial_j,\partial_i]u_s(x)
    &=
    \partial_j\partial_i u_s(x)-\partial_i\partial_j u_s(x)\\
    &=
    \tp{p_i(x)p_j(x)-p_j(x)p_i(x^{\oplus j})}
    \tp{u_s(x^{\oplus j})-u_s(x^{\oplus i})}\\
    &=
    \alpha_{ij}(x)
    \tp{\frac{\partial_j u_s(x)}{p_j(x)}
    -\frac{\partial_i u_s(x)}{p_i(x)}},
\end{align*}
where
\begin{align*}
    \alpha_{ij}(x):=p_i(x)p_j(x)-p_j(x)p_i(x^{\oplus j}).
\end{align*}
The square identity implies $\alpha_{ij}(x)=\alpha_{ji}(x)$.
Consequently,
\begin{align*}
    (L\partial_i-\partial_i L)u_s(x)
    =
    \sum_{j \neq i}[\partial_j,\partial_i]u_s(x)
    =
    \sum_{j \neq i}\alpha_{ij}(x)
    \tp{\frac{\partial_j u_s(x)}{p_j(x)}
    -\frac{\partial_i u_s(x)}{p_i(x)}}.
\end{align*}
The commutator contribution in the lower bound is therefore the following quadratic form in
$\nabla u_s(x)$:
\begin{align*}
    \sum_{i \in [n]}\partial_i u_s(x)(L\partial_i-\partial_i L)u_s(x)
    =
    \sum_{\substack{i,j \in [n]\\ i\neq j}}
    \alpha_{ij}(x)\partial_i u_s(x)
    \tp{\frac{\partial_j u_s(x)}{p_j(x)}
    -\frac{\partial_i u_s(x)}{p_i(x)}}.
\end{align*}
Equivalently, let
\begin{align*}
    g:=\tp{g_i}_{i\in[n]},\qquad
    g_i:=\partial_i u_s(x),\qquad
    r_i:=p_i(x),\qquad
    b_s:=1-\e^{-\rho s}.
\end{align*}
Since $p_i(x^{\oplus i})=1/r_i$, the lower bound above can be written as
\begin{align*}
    \abs{V_s(x)}(\partial_s+L)F_s(x)
    \ge g^\intercal A_s(x)g,
\end{align*}
where $A_s(x)$ is the symmetric matrix given by
\begin{align*}
    (A_s(x))_{ii}
    &=
    \frac{M\rho\e^{-\rho s}}{2}
    +\frac{2Mb_s-1}{r_i}
    -Mb_s\sum_{j\neq i}\frac{\alpha_{ij}(x)}{r_i},
    \\
    (A_s(x))_{ij}
    &=
    \frac{Mb_s}{2}\alpha_{ij}(x)
    \tp{\frac{1}{r_i}+\frac{1}{r_j}},
    \qquad i\neq j.
\end{align*}
Indeed, the commutator part can be symmetrized as
\begin{align*}
    \sum_{i \in [n]}g_i(L\partial_i-\partial_i L)u_s(x)
    =
    -\sum_{i<j}\alpha_{ij}(x)
    \tp{g_i-g_j}
    \tp{\frac{g_i}{r_i}-\frac{g_j}{r_j}}.
\end{align*}


\subsection{Zero-field Ising model}

\begin{claim}[Zero-field Ising model]
Use the $\frac12 x^\intercal Jx$ convention and assume
\begin{align*}
    \mu(x)=\frac{1}{Z}\exp\tp{\frac12\sum_{i,j=1}^n J_{ij}x_ix_j},
    \qquad J=J^\intercal,\qquad J_{ii}=0,
\end{align*}
where $x\in\{-1,+1\}^n$. Let
\begin{align*}
    B:=\max_i\sum_j |J_{ij}|.
\end{align*}
If $2B\e^{2B}<1$, then for the rates
\begin{align*}
    p_i(x)=\sqrt{\frac{\mu(x^{\oplus i})}{\mu(x)}}
\end{align*}
one may choose
\begin{align*}
    M\ge \frac{1}{2(1-2B\e^{2B})},
    \qquad
    \rho\ge \frac{2\e^B}{M},
\end{align*}
so that $A_s(x)\succeq0$ for all $s\ge0$ and all $x$. Consequently,
\begin{align*}
    I(P_t f)
    \le
    P_t\tp{\sqrt{I(f)^2+M\tp{1-\e^{-\rho t}}\abs{\nabla f}^2}}.
\end{align*}
\end{claim}

\begin{proof}
Let
\begin{align*}
    q:=\max_i\sum_{j\neq i}\tp{1-\e^{-2|J_{ij}|}},
    \qquad
    r_i(x):=p_i(x).
\end{align*}
Then
\begin{align*}
    r_i(x)=\exp\tp{-x_i\sum_jJ_{ij}x_j},
    \qquad
    \e^{-B}\le r_i(x),\, r_i(x)^{-1}\le \e^B.
\end{align*}
For $i\neq j$, put $\theta_{ij}=J_{ij}x_ix_j$ and
$H_i^{(j)}=\sum_{k\neq i,j}J_{ik}x_ix_k$. Then
\begin{align*}
    p_i(x^{\oplus j})=r_i(x)\e^{2\theta_{ij}},
    \qquad
    \alpha_{ij}(x)=r_i(x)r_j(x)\tp{1-\e^{2\theta_{ij}}},
\end{align*}
and the zero-field cancellation gives
\begin{align*}
    \frac{\alpha_{ij}(x)}{r_i(x)}
    =-2\e^{-H_j^{(i)}}\sinh(\theta_{ij}),
    \qquad
    \frac{\alpha_{ij}(x)}{r_j(x)}
    =-2\e^{-H_i^{(j)}}\sinh(\theta_{ij}).
\end{align*}
Thus
\begin{align}\label{eq:zero-field-alpha-bound}
    \abs{\frac{\alpha_{ij}(x)}{r_i(x)}},
    \abs{\frac{\alpha_{ij}(x)}{r_j(x)}}
    \le \e^B\tp{1-\e^{-2|J_{ij}|}}.
\end{align}

Let $b_s=1-\e^{-\rho s}$. We decompose
\begin{align*}
    A_s(x)=(1-b_s)E(x)+b_sC(x),
\end{align*}
where
\begin{align*}
    E_{ii}=\frac{M\rho}{2}-\frac1{r_i},
    \qquad E_{ij}=0\quad(i\neq j),
\end{align*}
and
\begin{align*}
    C_{ii}=\frac{2M-1}{r_i}-M\sum_{j\neq i}\frac{\alpha_{ij}}{r_i},
    \qquad
    C_{ij}=\frac{M}{2}\alpha_{ij}\tp{\frac1{r_i}+\frac1{r_j}}
    \quad(i\neq j).
\end{align*}
The matrix $E$ is positive semidefinite if $\rho\ge 2\e^B/M$. For $C$, Gershgorin and
\eqref{eq:zero-field-alpha-bound} give
\begin{align*}
    C_{ii}-\sum_{j\neq i}|C_{ij}|
    &\ge
    \frac{2M-1}{r_i}
    -M\sum_{j\neq i}\tp{
        \frac32\abs{\frac{\alpha_{ij}}{r_i}}
        +\frac12\abs{\frac{\alpha_{ij}}{r_j}}
    }\\
    &\ge (2M-1)\e^{-B}-2M\e^Bq.
\end{align*}
Hence $C(x)\succeq0$ for every $x$ if
\begin{align*}
    \e^{2B}q<1
    \qquad\text{and}\qquad
    M\ge \frac{1}{2(1-\e^{2B}q)}.
\end{align*}
Since $q\le 2B$, the claimed condition $2B\e^{2B}<1$ implies the previous display.
Equivalently,
\begin{align*}
    B<\frac12 W(1)\approx 0.28357.
\end{align*}
The stated choices of $M$ and $\rho$ therefore make both $E$ and $C$ positive
semidefinite. Since $A_s=(1-b_s)E+b_sC$ and $b_s\in[0,1]$, this proves
$A_s(x)\succeq0$ for all $s$ and $x$. The local Bobkov inequality follows from
the preceding semigroup argument.
\end{proof}

\section{Smoothing lemma for the zero-field Ising dynamics}

In this section we keep the zero-field Ising setting from the preceding
subsection and use the same square-root rates
\[
    p_i(x)=\sqrt{\frac{\mu(x^{\oplus i})}{\mu(x)}},
    \qquad
    P_t=\e^{tL}.
\]
Write
\[
    \delta_i f(x):=f(x^{\oplus i})-f(x).
\]
For \(q\ge1\), define the weighted edge norm
\[
    \norm{h}_{q,i}:=
    \tp{\E[\mu]{p_i(x)|h(x)|^q}}^{1/q},
\]
and define the \(L^1\)-influences by
\[
    I_i(f):=\norm{\delta_i f}_{1,i},
    \qquad
    W(f):=\sum_{i=1}^n I_i(f)^2.
\]
The Dirichlet form of \(L\) is
\[
    \mathcal E(g,g):=-\ip{g}{Lg}_{L^2(\mu)}
    =\frac12\sum_i\norm{\delta_i g}_{2,i}^2.
\]

We use two standard high-temperature inputs.  First, assume that the dynamics
with generator \(L\), and the auxiliary dynamics obtained by flipping or
freezing one coordinate, satisfy a uniform log-Sobolev inequality: there is
\(\alpha>0\), independent of \(n\), such that
\[
    \Ent[\nu]{u^2}\le \frac{2}{\alpha}\mathcal E_\nu(u,u).
\]
Equivalently, their semigroups satisfy hypercontractivity,
\[
    \norm{Q_t h}_{q,\nu}\le \norm{h}_{p,\nu},
    \qquad
    q-1=(p-1)\e^{2\alpha t}.
\]
Second, define
\[
    a_{ij}:=\sup_x |p_j(x^{\oplus i})-p_j(x)|,\qquad a_{ii}:=0,
    \qquad \mathsf A=(a_{ij})_{i,j\in[n]},
\]
and assume
\[
    \norm{\mathsf A}_{\ell^2\to\ell^2}\le\eta
\]
for \(\eta>0\) sufficiently small.  In the zero-field Ising model,
\[
    p_i(x)=\exp\tp{-x_i\sum_{j\ne i}J_{ij}x_j},
    \qquad
    \e^{-B}\le p_i(x)\le \e^B,
\]
and, for \(i\ne j\),
\[
    p_j(x^{\oplus i})=p_j(x)\e^{2J_{ij}x_ix_j}.
\]
Therefore
\[
    a_{ij}\le \e^B(\e^{2|J_{ij}|}-1),
    \qquad
    \norm{\mathsf A}_{\ell^2\to\ell^2}\le \e^B(\e^{2B}-1).
\]
Thus the commutator assumption follows, for instance, from \(B\) sufficiently
small.

\begin{lemma}[Stability of \(L^1\)-influences]\label{lem:smooth-l1-stability}
If \(\norm{\mathsf A}_{\ell^2\to\ell^2}\le 2\e^{-4B}\), then
\[
    W(P_t f)\le \e^{12B} W(f)
\]
for all \(t\ge0\).
\end{lemma}

\begin{proof}
In this proof,
\[
    \norm{h}_{L^1(\mu)}:=\int |h|\,d\mu=\E[\mu]{|h|}
\]
denotes the unweighted \(L^1(\mu)\)-norm.  Since
\(\e^{-B}\le p_i\le\e^B\),
\[
    \e^{-B}\norm{h}_{L^1(\mu)}
    \le
    \norm{h}_{1,i}
    \le
    \e^B\norm{h}_{L^1(\mu)}.
\]
Thus it is enough to control the unweighted edge \(L^1\)-norms, keeping track
of these factors at the end.

Let \(g_t:=P_t f\), and write \(p_j^{\oplus i}(x):=p_j(x^{\oplus i})\).
A direct computation gives
\[
    \delta_i Lg
    =
    -(p_i+p_i^{\oplus i})\delta_i g
    +\sum_{j\ne i}p_j^{\oplus i}\delta_j\delta_i g
    +\sum_{j\ne i}(p_j^{\oplus i}-p_j)\delta_jg.
\]
Define
\[
    L^{(i)}h:=\sum_{j\ne i}p_j(x^{\oplus i})\delta_jh,
    \qquad
    r_i:=p_i+p_i^{\oplus i}.
\]
Since \(p_i^{\oplus i}=p_i^{-1}\), \(r_i\ge2\), and hence
\[
    \partial_t\delta_i g_t
    =
    (L^{(i)}-r_i)\delta_i g_t
    +\sum_{j\ne i}(p_j^{\oplus i}-p_j)\delta_jg_t.
\]

Let \(S_t^{(i)}\) be the semigroup generated by \(L^{(i)}-r_i\).  By
positivity and Feynman--Kac,
\[
    |S_t^{(i)}h|\le \e^{-2t}Q_t^{(i)}|h|,
\]
where \(Q_t^{(i)}\) is generated by \(L^{(i)}\).  This Markov semigroup is
reversible with respect to \(\nu_i(x):=\mu(x^{\oplus i})\), and
\[
    \frac{d\nu_i}{d\mu}(x)=p_i(x)^2
\]
lies in \([\e^{-2B},\e^{2B}]\).  Therefore, writing
\(w_i=d\mu/d\nu_i=p_i^{-2}\),
\[
    \int Q_t^{(i)}|h|\,d\mu
    =
    \int Q_t^{(i)}|h|\, w_i\,d\nu_i
    \le
    \e^{2B}\int Q_t^{(i)}|h|\,d\nu_i
    =
    \e^{2B}\int |h|\,d\nu_i
    \le
    \e^{4B}\norm{h}_{L^1(\mu)}.
\]
Thus, by Feynman--Kac, it holds
\[
    \norm{S_t^{(i)}h}_{L^1(\mu)}
    \le \e^{4B}\e^{-2t}\norm{h}_{L^1(\mu)}.
\]
Duhamel's formula gives
\[
    \delta_i g_t
    =
    S_t^{(i)}\delta_i f
    +
    \int_0^t
    S_{t-s}^{(i)}
    \sum_{j\ne i}(p_j^{\oplus i}-p_j)\delta_jg_s\,ds,
\]
and hence
\[
    \norm{\delta_i g_t}_{L^1(\mu)}
    \le
    \e^{4B}\e^{-2t}\norm{\delta_i f}_{L^1(\mu)}
    +
    \e^{4B}\sum_{j\ne i}a_{ij}
    \int_0^t\e^{-2(t-s)}\norm{\delta_jg_s}_{L^1(\mu)}\,ds.
\]
Set
\[
    Q(t):=\tp{\sum_i\norm{\delta_i g_t}_{L^1(\mu)}^2}^{1/2}.
\]
Equivalently, let
\[
    q_i(t):=\norm{\delta_i g_t}_{L^1(\mu)},
    \qquad
    q(t):=(q_i(t))_{i\in[n]}.
\]
The preceding estimate says componentwise that
\[
    q(t)
    \le
    \e^{4B}\e^{-2t}q(0)
    +
    \e^{4B}\int_0^t\e^{-2(t-s)}\mathsf A q(s)\,ds,
\]
where we use \(a_{ii}=0\) and \(a_{ij}\ge0\).  Taking the \(\ell^2\)-norm
and applying Minkowski's integral inequality gives
\[
    \norm{q(t)}_2
    \le
    \e^{4B}\e^{-2t}\norm{q(0)}_2
    +
    \e^{4B}\int_0^t\e^{-2(t-s)}
    \norm{\mathsf A q(s)}_2\,ds.
\]
Since \(Q(t)=\norm{q(t)}_2\) and
\(\norm{\mathsf A q(s)}_2\le\norm{\mathsf A}_{2\to2}Q(s)\), this yields
\[
    Q(t)
    \le
    \e^{4B}\e^{-2t}Q(0)
    +
    \e^{4B}\norm{\mathsf A}_{2\to2}
    \int_0^t\e^{-2(t-s)}Q(s)\,ds.
\]
Multiplying by \(\e^{2t}\) and applying Gronwall,
\[
    Q(t)
    \le
    \e^{4B}
    \e^{-(2-\e^{4B}\norm{\mathsf A}_{2\to2})t}Q(0).
\]
The assumed bound \(\norm{\mathsf A}_{2\to2}\le2\e^{-4B}\) gives
\[
    Q(t)\le \e^{4B}Q(0).
\]
Finally,
\[
    W(P_t f)
    =
    \sum_i\norm{\delta_i g_t}_{1,i}^2
    \le
    \e^{2B}Q(t)^2
    \le
    \e^{10B}Q(0)^2
    \le
    \e^{12B}W(f),
\]
where the last inequality uses
\(\norm{\delta_i f}_{L^1(\mu)}\le\e^B\norm{\delta_i f}_{1,i}\).
\end{proof}

\paragraph{Constants used in the smoothing argument.}
The constants \(M\) and \(\rho\) in the local Bobkov inequality above are
chosen from the pointwise matrix condition \(A_s(x)\succeq0\).  They are not
the constants used to run the smoothing argument below.  From the
commutator estimate we get the gradient-decay margin
\[
    \lambda:=2-\e^{4B}\norm{\mathsf A}_{2\to2}.
\]
The log-Sobolev constant is the constant \(\alpha\) in the
hypercontractive relation \(q-1=(p-1)\e^{2\alpha t}\).  It enters the
remaining constants only through
\[
    \beta:=\tanh\alpha,
\]
because the interpolation exponent below is \(\tanh(\alpha t)\).  For the
two results below assume \(\lambda>0\), and set
\[
    C_{\mathrm{int}}
    :=
    \frac{\e^B}{2\lambda}
    +2+\frac1\beta
    +\frac{1}{2\lambda}\max\left\{1,\frac{1}{2\beta}\right\},
    \qquad
    C_{\mathrm{Tal}}:=\e^{11B}C_{\mathrm{int}}.
\]

\begin{lemma}[Talagrand-type variance inequality]\label{lem:smooth-talagrand}
Under the log-Sobolev assumption with constant \(\alpha\) and \(\lambda>0\),
every \(g:X\to\R\) satisfies
\[
    \Var[\mu]{g}
    \le
    C_{\mathrm{Tal}}\sum_i
    \frac{\norm{\delta_i g}_{2,i}^2}
    {1+\log_+\tp{\norm{\delta_i g}_{2,i}/\norm{\delta_i g}_{1,i}}}.
\]
Here \(C_{\mathrm{Tal}}=\e^{11B}C_{\mathrm{int}}\), with
\(C_{\mathrm{int}}\) defined above.  Thus \(C_{\mathrm{Tal}}\) depends on
\(\alpha\) only through \(\beta=\tanh\alpha\), and
\[
    C_{\mathrm{Tal}}\asymp_{B,\lambda}1+\frac1{\tanh\alpha}.
\]
The \(i\)-term is interpreted as \(0\) if \(\delta_i g=0\).
\end{lemma}

\begin{remark}
    This should be essentially the same as the one in Koehler's paper. Though, the degree might be unbounded in the underlying graph.
\end{remark}
\begin{proof}
We first prove that, for \(p_t:=1+\e^{-2\alpha t}\),
\[
    \tp{\sum_i\norm{\delta_iP_tg}_{2,i}^2}^{1/2}
    \le
    \e^{11B/2}\e^{-\lambda t}
    \tp{\sum_i\norm{\delta_i g}_{p_t,i}^2}^{1/2}.
\]
It suffices to prove the corresponding unweighted estimate.  For
\(1\le p\le q\le2\), the comparison between \(\mu\) and
\(\nu_i(x)=\mu(x^{\oplus i})\), Feynman--Kac, and hypercontractivity give
\[
    \norm{S_s^{(i)}h}_{q,\mu}
    \le
    \e^{4B}\e^{-2s}\norm{h}_{p,\mu},
    \qquad q-1=(p-1)\e^{2\alpha s}.
\]

Let \(M_{ij}\) denote multiplication by \(p_j^{\oplus i}-p_j\), with
\(M_{ii}=0\).  The equation from \cref{lem:smooth-l1-stability} gives
\[
    \delta_iP_tg
    =
    S_t^{(i)}\delta_i g
    +
    \int_0^tS_{t-s}^{(i)}\sum_jM_{ij}\delta_jP_sg\,ds.
\]
Iterating yields a sum over paths \(i=i_0,i_1,\ldots,i_m\) and time
increments \(\tau_0,\ldots,\tau_m\ge0\) with
\(\tau_0+\cdots+\tau_m=t\):
\[
    S_{\tau_0}^{(i_0)}M_{i_0i_1}
    S_{\tau_1}^{(i_1)}M_{i_1i_2}\cdots
    M_{i_{m-1}i_m}S_{\tau_m}^{(i_m)}\delta_{i_m}g.
\]
Choosing the intermediate exponents so that, after elapsed time \(\sigma\),
the exponent is \(1+\e^{-2\alpha(t-\sigma)}\), each path term is bounded in
\(L^2(\mu)\) by
\[
    \e^{4B(m+1)}\e^{-2t}
    \tp{\prod_{\ell=0}^{m-1}a_{i_\ell i_{\ell+1}}}
    \norm{\delta_{i_m}g}_{p_t}.
\]
After integrating over the \(m\)-simplex,
\[
    \norm{\delta_iP_tg}_2
    \le
    \e^{4B}\e^{-2t}
    \sum_{m\ge0}\frac{(\e^{4B}t)^m}{m!}
    \sum_j(\mathsf A^m)_{ij}\norm{\delta_jg}_{p_t}.
\]
Taking the \(\ell^2\)-norm in \(i\),
\[
    \tp{\sum_i\norm{\delta_iP_tg}_2^2}^{1/2}
    \le
    \e^{4B}\e^{-\lambda t}
    \tp{\sum_i\norm{\delta_i g}_{p_t}^2}^{1/2}.
\]
Converting the left side to weighted \(L^2\)-edge norms costs at most
\(\e^{B/2}\), and converting the right side from unweighted \(L^{p_t}\) to
weighted \(L^{p_t}\)-edge norms costs at most \(\e^B\).  This proves the
displayed weighted gradient estimate.

By reversibility and irreducibility,
\[
    \Var[\mu]{g}
    =
    \int_0^\infty\sum_i\norm{\delta_iP_tg}_{2,i}^2\,dt.
\]
The gradient estimate gives
\[
    \Var[\mu]{g}
    \le
    \e^{11B}\int_0^\infty\e^{-2\lambda t}
    \sum_i\norm{\delta_i g}_{p_t,i}^2\,dt.
\]

Fix \(i\), and set
\[
    a_i:=\norm{\delta_i g}_{1,i},\qquad b_i:=\norm{\delta_i g}_{2,i}.
\]
If \(b_i=0\), there is nothing to prove; otherwise \(a_i>0\).  Interpolation
gives, for \(1\le p\le2\),
\[
    \norm{\delta_i g}_{p,i}
    \le
    a_i^{(2-p)/p}b_i^{2(p-1)/p}.
\]
Since \((2-p_t)/p_t=\tanh(\alpha t)\),
\[
    \norm{\delta_i g}_{p_t,i}^2
    \le
    b_i^2\tp{\frac{a_i}{b_i}}^{2\tanh(\alpha t)}.
\]
Write \(R_i:=b_i/a_i\).  If \(R_i<1\), then
\[
    a_i\le \tp{\int p_i\,d\mu}^{1/2}b_i\le \e^{B/2}b_i,
\]
so \(R_i\ge\e^{-B/2}\) and
\[
    \int_0^\infty\e^{-2\lambda t}\norm{\delta_i g}_{p_t,i}^2\,dt
    \le
    \frac{\e^B}{2\lambda}b_i^2.
\]
If \(R_i\ge1\), put \(L_i:=\log R_i\).  Since
\(\tanh(\alpha t)\ge\beta\min\{t,1\}\),
\[
    \int_0^\infty\e^{-2\lambda t}R_i^{-2\tanh(\alpha t)}\,dt
    \le
    \frac{2+1/\beta}{1+L_i}
    +
    \frac{1}{2\lambda}\max\left\{1,\frac{1}{2\beta}\right\}
    \frac{1}{1+L_i}.
\]
Combining the two cases gives
\[
    \int_0^\infty\e^{-2\lambda t}\norm{\delta_i g}_{p_t,i}^2\,dt
    \le
    C_{\mathrm{int}}
    \frac{b_i^2}{1+\log_+R_i}.
\]
Summing over \(i\) proves the lemma.
\end{proof}

Set
\[
    D:=3C_{\mathrm{Tal}},
    \qquad
    \gamma:=\frac{1}{D(2+\log D)},
    \qquad
    \rho_{\mathrm{sm}}:=\frac{\gamma}{2}.
\]
Thus the final smoothing rate is
\[
    \rho_{\mathrm{sm}}
    =
    \frac{1}{2D(2+\log D)},
    \qquad
    D=3\e^{11B}C_{\mathrm{int}}.
\]
In particular, the log-Sobolev constant affects the final rate through
\(\beta=\tanh\alpha\): a smaller \(\alpha\) makes \(\beta\) smaller, increases
\(C_{\mathrm{int}}\) and \(C_{\mathrm{Tal}}\), and therefore decreases
\(\rho_{\mathrm{sm}}\).  For fixed \(B\) and \(\lambda\),
\[
    \rho_{\mathrm{sm}}
    \asymp_{B,\lambda}
    \frac{\tanh\alpha}{1+\log(1+1/\tanh\alpha)}.
\]

\begin{theorem}[Smoothing lemma]\label{thm:smoothing-ising}
Under the log-Sobolev assumption with constant \(\alpha\) and \(\lambda>0\),
every Boolean \(f:X\to\{-1,+1\}\) satisfies
\[
    \Var[\mu]{P_t f}
    \le
    \e^{12B}\,W(f)^{\theta_t}\Var[\mu]{f}^{1-\theta_t},
    \qquad
    \theta_t=\tanh(\rho_{\mathrm{sm}}t).
\]
Here
\[
    \rho_{\mathrm{sm}}
    =
    \frac{1}{2D(2+\log D)},
    \qquad
    D=3\e^{11B}C_{\mathrm{int}},
\]
so, for fixed \(B\) and \(\lambda\),
\[
    \rho_{\mathrm{sm}}
    \asymp_{B,\lambda}
    \frac{\tanh\alpha}{1+\log(1+1/\tanh\alpha)}.
\]
\end{theorem}

\begin{proof}
If \(\Var[\mu]{f}=0\), then \(f\) is constant because \(\mu\) has full support.
If \(W(f)=0\), then \(\delta_i f=0\) on every edge; since the hypercube is
connected, \(f\) is again constant.  We may therefore assume
\[
    V:=\Var[\mu]{f}>0,
    \qquad W(f)>0.
\]

Let
\[
    F(s):=\Var[\mu]{P_s f},\qquad
    \mathcal B(s):=\sum_i\norm{\delta_iP_s f}_{2,i}^2,\qquad
    \mathcal A(s):=\sum_i\norm{\delta_iP_s f}_{1,i}^2.
\]
By reversibility, \(F'(s)=-\mathcal B(s)\).  By
\cref{lem:smooth-l1-stability},
\[
    \mathcal A(s)=W(P_s f)\le \e^{12B} W(f).
\]

Apply \cref{lem:smooth-talagrand} to \(g=P_s f\).  With
\[
    a_i:=\norm{\delta_iP_s f}_{1,i},
    \qquad
    b_i:=\norm{\delta_iP_s f}_{2,i},
\]
we get
\[
    F(s)\le
    C_{\mathrm{Tal}}\sum_i\frac{b_i^2}{1+\log_+(b_i/a_i)}.
\]
Then \(\mathcal A(s)=\sum_i a_i^2\) and
\(\mathcal B(s)=\sum_i b_i^2\).  Indices with \(a_i=0\) also have
\(b_i=0\).  If \(\mathcal A(s)=0\), then \(F(s)=0\), and the claim is
trivial.  Otherwise,
\[
    \sum_i\frac{b_i^2}{1+\log_+(b_i/a_i)}
    \le
    3\frac{\mathcal B(s)}
    {1+\log_+\sqrt{\mathcal B(s)/\mathcal A(s)}}.
\]
Indeed, if \(\mathcal B(s)\le\e^4\mathcal A(s)\), then the right side is at
least \(\mathcal B(s)\).  Otherwise split at
\((\mathcal B(s)/\mathcal A(s))^{1/4}\):
The large-ratio part can be bounded by
\begin{align*}
    \sum_{b_i / a_i \ge \tp{\+B(s)/\+A(s)}^{1/4}} \frac{b_i^2}{1+\log_+(b_i/a_i)} \le \sum_{i} \frac{2 b_i^2}{1+\log_+\tp{\sqrt{\+B(s)/\+A(s)}}} \le \frac{2\+B(s)}{1+\log_+(\sqrt{\+B(s)/A(s)})}.
\end{align*}
The
small-ratio part costs at most
\begin{align*}
\sum_{b_i / a_i \le \tp{\+B(s)/\+A(s)}^{1/4}} \frac{(b_i/a_i)^2 a_i^2}{1+\log_+(b_i/a_i)} \le \sqrt{\frac{\+B(s)}{\+A(s)}}\sum_{i} a_i^2 = \sqrt{\+A(s) \+B(s)} \le \frac{\+B(s)}{1 + \log_+ \sqrt{\+B(s)/\+A(s)}}.
\end{align*}


With \(K:=\e^{12B}W(f)\), the preceding bounds give
\[
    F(s)
    \le
    D\frac{\mathcal B(s)}
    {1+\log_+\sqrt{\mathcal B(s)/K}}.
\]
The elementary implication
\[
    y\le D\frac{x}{1+\log_+\sqrt{x/K}}
    \quad\Longrightarrow\quad
    x\ge \gamma\,y\tp{1+\log_+\frac{y}{K}}
\]
follows from \(y\le Dx\) and
\[
    1+\log_+\frac{y}{K}
    \le
    (2+\log D)\tp{1+\log_+\sqrt{x/K}}.
\]
Applying this with \(x=\mathcal B(s)\) and \(y=F(s)\), we obtain
\[
    F'(s)
    \le
    -\gamma F(s)\tp{1+\log_+\frac{F(s)}{K}}.
\]

If \(K\ge V\), then
\[
    F(t)\le V\le \e^{12B}W(f)^{\theta_t}V^{1-\theta_t}.
\]
Assume \(K<V\), and let \(\tau:=\inf\{s:F(s)\le K\}\).  On
\([0,\tau)\), set \(u(s):=\log(F(s)/K)\).  Then
\[
    u'(s)=\frac{F'(s)}{F(s)}\le -\gamma(1+u(s))\le -\gamma u(s),
\]
so
\[
    F(s)\le K^{1-\e^{-\gamma s}}V^{\e^{-\gamma s}}
    \qquad (0\le s<\tau).
\]
If \(\tau\le t\), monotonicity gives \(F(t)\le K\), which satisfies the same
bound.  Thus, for all \(t\ge0\),
\[
    F(t)\le K^{1-\e^{-\gamma t}}V^{\e^{-\gamma t}}.
\]
Since \(\theta_t=\tanh(\rho_{\mathrm{sm}}t)=\tanh(\gamma t/2)
\le1-\e^{-\gamma t}\) and \(K<V\),
\[
    F(t)
    \le
    K^{\theta_t}V^{1-\theta_t}
    \le
    \e^{12B}W(f)^{\theta_t}\Var[\mu]{f}^{1-\theta_t}.
\]
\end{proof}

\section{Eldan--Gross inequality for high-temperature Ising}

We keep the zero-field Ising notation, square-root rates, and definitions of
\(\delta_i,\partial_i,|\nabla f|,\|\cdot\|_{q,i},I_i(f)\), and \(W(f)\) from
the previous section.  Let \(M,\rho_{\mathrm B}\) denote the constants in the
local Bobkov inequality, and use the smoothing estimate above with constants
\[
    C_0=\e^{12B},\qquad C_1=1,\qquad
    \rho_{\mathrm{sm}}=\frac{1}{2D(2+\log D)}.
\]
We also use the hypercontractivity corresponding to the log-Sobolev constant
\(\alpha\), and write \(\lambda_{\mathrm{gap}}\) for the spectral gap:
\[
    \Var[\mu]{P_t g}
    \le
    \e^{-2\lambda_{\mathrm{gap}}t}\Var[\mu]{g}.
\]
All these constants are dimension-free in the high-temperature regime.

We first record the variance-drop consequence of local Bobkov.

\begin{lemma}[Variance drop from local Bobkov]\label{lem:variance-drop-bobkov}
There is a constant \(C_{\mathrm B}>0\), depending only on \(M\), such that
every Boolean \(f:X\to\{-1,+1\}\) satisfies
\[
    \Var[\mu]{f}-\Var[\mu]{P_t f}
    \le
    C_{\mathrm B}
    \sqrt{1-\e^{-\rho_{\mathrm B}t}}\,
    \E[\mu]{|\nabla f|}.
\]
Consequently,
\[
    \Var[\mu]{f}-\Var[\mu]{P_t f}
    \le
    C_{\mathrm B}
    \sqrt{\rho_{\mathrm B}t}\,
    \E[\mu]{|\nabla f|}.
\]
\end{lemma}

\begin{proof}
Put \(h=(1+f)/2\).  Then \(I(h)=0\), \(P_th=(1+P_tf)/2\), and
\(|\nabla h|=|\nabla f|/2\).  Local Bobkov gives
\[
    I\left(\frac{1+P_t f}{2}\right)
    \le
    \frac{\sqrt M}{2}
    \sqrt{1-\e^{-\rho_{\mathrm B}t}}\,
    P_t|\nabla f|.
\]
Using the universal bound
\[
    I\left(\frac{1+s}{2}\right)\ge c_I(1-s^2),
    \qquad -1\le s\le1,
\]
and stationarity of \(P_t\), we obtain
\[
    c_I\E[\mu]{1-(P_t f)^2}
    \le
    \frac{\sqrt M}{2}
    \sqrt{1-\e^{-\rho_{\mathrm B}t}}\,
    \E[\mu]{|\nabla f|}.
\]
Since \(f^2=1\) and \(\E[\mu]{P_tf}=\E[\mu]{f}\),
\[
    \E[\mu]{1-(P_t f)^2}
    =
    \Var[\mu]{f}-\Var[\mu]{P_t f}.
\]
This proves the first bound with \(C_{\mathrm B}:=\sqrt M/(2c_I)\).  The
second follows from \(1-\e^{-u}\le u\).
\end{proof}

Next we prove a preliminary lower bound which captures the
\(\log(e/\Var f)\) contribution.

\begin{lemma}[Large-variance and hypercontractive lower bound]
\label{lem:variance-log-lower}
There is a constant \(c_1>0\), depending only on the high-temperature
constants, such that every Boolean \(f:X\to\{-1,+1\}\) satisfies
\[
    \E[\mu]{|\nabla f|}
    \ge
    c_1\Var[\mu]{f}
    \sqrt{\log\left(\frac{e}{\Var[\mu]{f}}\right)}.
\]
\end{lemma}

\begin{proof}
Let \(V:=\Var[\mu]{f}\) and \(G:=\E[\mu]{|\nabla f|}\).  The case \(V=0\) is
trivial.

If \(V\ge\e^{-2}\), take
\[
    t_0:=\frac{\log 2}{2\lambda_{\mathrm{gap}}}.
\]
The spectral-gap estimate gives \(\Var[\mu]{P_{t_0}f}\le V/2\).  Hence
\cref{lem:variance-drop-bobkov} implies
\[
    \frac V2
    \le C_{\mathrm B}\sqrt{1-\e^{-\rho_{\mathrm B}t_0}}\,G,
\]
so \(G\ge cV\).  Since \(\log(e/V)\le3\) in this case, the desired bound
follows.

Assume now \(0<V<\e^{-2}\), and set \(g:=f-\E[\mu]{f}\).  For Boolean \(f\),
\(\norm{g}_2^2=\norm{g}_1=V\).  With
\[
    p_t:=1+\e^{-2\alpha t},
    \qquad
    \eta_t:=\frac{2-p_t}{p_t}=\tanh(\alpha t),
\]
hypercontractivity and interpolation give
\[
    \Var[\mu]{P_t f}
    =\norm{P_tg}_2^2
    \le \norm{g}_{p_t}^2
    \le V^{1+\eta_t}.
\]
Let \(L_V:=\log(1/V)\ge2\) and \(t_1:=1/(\alpha L_V)\).  Then
\(\eta_{t_1}=\tanh(1/L_V)\ge c/L_V\), hence
\(\Var[\mu]{P_{t_1}f}\le e^{-c}V\).  Applying
\cref{lem:variance-drop-bobkov} at time \(t_1\),
\[
    cV
    \le C_{\mathrm B}\sqrt{\rho_{\mathrm B}t_1}\,G,
    \qquad\text{so}\qquad
    G\ge c_1V\sqrt{L_V}.
\]
Since \(L_V\asymp\log(e/V)\) for \(V<\e^{-2}\), this proves the claim after
decreasing \(c_1\).
\end{proof}

We next extract the \(W(f)\)-dependent logarithm from the smoothing lemma.

\begin{lemma}[Smoothing lower bound]
\label{lem:smoothing-log-lower}
There are constants \(c_2>0\) and \(C_2\ge1\), depending only on the
high-temperature constants, such that every Boolean \(f:X\to\{-1,+1\}\)
satisfies
\[
    \E[\mu]{|\nabla f|}
    \ge
    c_2
    \Var[\mu]{f}
    \sqrt{
        \log_+\left(
            \frac{\Var[\mu]{f}}{C_2W(f)}
        \right)
    }.
\]
\end{lemma}

\begin{proof}
Let \(V:=\Var[\mu]{f}\) and \(G:=\E[\mu]{|\nabla f|}\).  If \(V=0\), or if
\(W(f)=0\) (which forces \(f\) to be constant), there is nothing to prove.

Set
\[
    A_0:=\max\{1,2\log(2C_0)\},
    \qquad
    C_2:=e^{A_0}C_1.
\]
If \(V\le C_2W(f)\), the right-hand side is zero.  Otherwise put
\[
    R:=\log\left(\frac{V}{C_1W(f)}\right)
    =
    \log\left(\frac{V}{C_2W(f)}\right)+A_0
    \ge A_0,
    \qquad
    t_2:=\frac{A_0}{\rho_{\mathrm{sm}}R}.
\]
Since \(A_0/R\le1\),
\[
    \theta_{t_2}=\tanh(A_0/R)\ge \frac{A_0}{2R}.
\]
The smoothing lemma gives
\[
    \Var[\mu]{P_{t_2}f}
    \le C_0V\e^{-\theta_{t_2}R}
    \le \frac V2,
\]
by the choice of \(A_0\).  The variance-drop lemma then yields
\[
    \frac V2
    \le C_{\mathrm B}\sqrt{\rho_{\mathrm B}t_2}\,G,
    \qquad\text{hence}\qquad
    G\ge c_2V\sqrt{R}.
\]
Finally \(R\ge\log(V/(C_2W(f)))\), and the claim follows after adjusting
\(c_2\).
\end{proof}

We now combine the two estimates.

\begin{theorem}[Eldan--Gross inequality for high-temperature Ising]
\label{thm:eg-high-temp-ising}
Under the assumptions above, there exists a constant
\(c_{\mathrm{HT}}>0\), depending only on the high-temperature constants and
not on \(n\), such that every Boolean \(f:X\to\{-1,+1\}\) satisfies
\[
    \E[\mu]{|\nabla f|}
    \ge
    c_{\mathrm{HT}}\,
    \Var[\mu]{f}
    \sqrt{
        \log\left(
            1+\frac{e}{W(f)}
        \right)
    }.
\]
\end{theorem}

\begin{proof}
Let \(V:=\Var[\mu]{f}\) and \(G:=\E[\mu]{|\nabla f|}\).  If \(V=0\), or if
\(W(f)=0\) (which forces \(f\) to be constant), the result is trivial.  The
two previous lemmas give, with \(c_3:=2^{-1/2}\min\{c_1,c_2\}\),
\[
    G
    \ge
    c_3V
    \sqrt{
        \log(e/V)
        +
        \log_+\left(\frac{V}{C_2W(f)}\right)
    }.
\]
It remains to use the elementary comparison
\[
    \log(e/V)
    +
    \log_+\left(\frac{V}{C_2W(f)}\right)
    \ge
    c(C_2)
    \log\left(1+\frac{e}{W(f)}\right),
\]
valid for \(0<V\le1\).  If \(V\ge C_2W(f)\), the left side equals
\[
    \log\left(\frac{e}{C_2W(f)}\right),
\]
which is comparable to \(\log(1+e/W(f))\) because \(W(f)\le1/C_2\).
If \(V<C_2W(f)\), the second term vanishes and
\[
    \log\left(1+\frac{e}{W(f)}\right)
    \le
    \log\left(1+\frac{eC_2}{V}\right)
    \lesssim_{C_2}
    \log(e/V).
\]
Thus
\[
    G
    \ge
    c_{\mathrm{HT}}V
    \sqrt{
        \log\left(1+\frac{e}{W(f)}\right)
    },
\]
as claimed.
\end{proof}

\begin{remark}[Dependence of the constant]
The constant \(c_{\mathrm{HT}}\) depends on the local Bobkov constants
\(M,\rho_{\mathrm B}\), the log-Sobolev constant \(\alpha\), the spectral
gap \(\lambda_{\mathrm{gap}}\), and the smoothing constants
\(C_0,C_1,\rho_{\mathrm{sm}}\). In the fixed high-temperature regime all of
these are dimension-free constants depending only on the interaction
strength parameter \(B\) and the commutator margin \(\lambda\). Hence
\[
    c_{\mathrm{HT}}=\Theta_{B,\lambda,\alpha}(1).
\]
In particular, for fixed high-temperature parameters,
\[
    \E[\mu]{|\nabla f|}
    \gtrsim
    \Var[\mu]{f}
    \sqrt{
        \log\left(
            1+\frac{e}{W(f)}
        \right)
    },
\]
with a constant independent of \(n\).
\end{remark}
\end{document}
